% !TeX root = ../main.tex
\section{Motivation}
\label{sec:challenges}
Figure~\ref{fig:insights} examines cold start, repeated reads, and changes in the active expert set.

\subsection{Cold Start}
\label{sec:challenge_initial}
Cold start concerns expert availability when the first input arrives.
Initial residents determine which selected experts require HDD reads.
Expert demand varies within and across layers~\cite{routingconsistency}.
A fixed slot count can cover concentrated demand in one layer and leave more accesses uncovered in another.
Figure~\ref{fig:insights}(a) compares Separate, Streaming, Fiddler~\cite{fiddler}, and MECC using TTFT normalized to Separate.
Figure~\ref{fig:insights}(b) expands one cold expert access.
HDD reading, host staging, and H2D precede expert execution.
The first token waits for these transfers.
Startup placement therefore involves both per-layer capacities and the expert IDs assigned to those slots.

\subsection{Repeated Reads}
\label{sec:challenge_reuse}
An expert evicted from HBM may be requested again during generation.
A CPU copy serves that request through H2D.
Figure~\ref{fig:insights}(c) compares TPOT normalized to Separate across the same four methods.
Independent replacement can retain the same experts in both tiers.
In Figure~\ref{fig:insights}(d), four slots hold only $b$ and $c$.
A later request for $a$ misses both caches and reloads its weights from HDD.
Repeated transfers add weight-delivery time to otherwise identical expert computation.
Duplicate residents occupy space that could retain other reusable experts.

\subsection{Dynamic Cache Adaptation}
\label{sec:challenge_shift}
Task changes alter which experts receive requests.
An initial ranking reflects the inputs analyzed before deployment.
Figure~\ref{fig:insights}(e) compares fixed and online CPU retention.
Both use the same initial residents, HBM LRU, and slot budgets.
Figure~\ref{fig:insights}(f) follows requests $u,v,w,u$.
The request for $w$ replaces $u$ in HBM.
CPU memory still retains the earlier task's $a,b,c,d$.
The final $u$ therefore misses both caches and reads its weights again.
The four CPU slots could retain all three task-B experts.
Their fixed membership leaves this capacity unavailable to the new working set.
Membership updates redirect these slots toward the serving task's experts.
